\chapter[Hashes SHA-256 das partições congeladas]{Hashes SHA-256 das partições congeladas}
\label{ap:sha256}

Para garantir a reprodutibilidade dos experimentos, as partições de treino,
desenvolvimento e teste foram congeladas após sua geração (semente $42$,
estratificação multirrótulo em nível de documento). O Quadro~\ref{tab:sha256}
apresenta os \emph{hashes} SHA-256 dos arquivos, gerados por
\texttt{scripts/splitting/make\_splits.py} e versionados em
\texttt{data/splits/SHA256SUMS.txt}. Qualquer alteração nos arquivos altera o
\emph{hash}, permitindo a verificação independente da integridade das partições.

\begin{quadro}[htbp]
  \centering
  \caption{SHA-256 das partições congeladas do SemClinBr.}
  \label{tab:sha256}
  {\footnotesize
  \begin{tabular}{ll}
  \toprule
  Arquivo & SHA-256 \\
  \midrule
  \texttt{train.jsonl} & \texttt{3d7a2a0730f8931a405449a4062f529700a632815767e25336c422e101933775} \\
  \texttt{dev.jsonl}   & \texttt{20ce38a9807d7488002fc7de865096ca6357c64021e44f6dd78f92c3346d7837} \\
  \texttt{test.jsonl}  & \texttt{597af7fb88c2bedeec304ea7c92afbe3f6851e913a10d643f2ce34511e77daed} \\
  \bottomrule
  \end{tabular}}\\[4pt]
  {\footnotesize Fonte: elaborado pelo autor (\texttt{data/splits/SHA256SUMS.txt}).}
\end{quadro}

As partições contêm, respectivamente, $800$, $100$ e $100$ documentos. A
distribuição das relações por partição --- $9\,105$ no treino, $1\,123$ no
desenvolvimento e $1\,125$ no teste --- preserva a proporção entre as classes
\texttt{associated\_with} e \texttt{negation\_of} observada no corpus completo.
